%% Part 3 — normal forms from subgroup chains (paper §3).
\PaperPart
\section{Normal forms from subgroup chains}

\begin{frame}{Transversals are digits \InPaper}
\begin{columns}[c]
\begin{column}{0.50\textwidth}
$H \le G$, a \alert{right transversal} picks one $\bar g$ per coset $Hg$:
\[ g = h\cdot\bar g \quad\text{uniquely},\qquad G \cong H\times (H\backslash G)\ \text{(as sets)}. \]
Fix a chain with a transversal $T_k$ for each step:
\[ 1 = G_0 \le G_1 \le \dots \le G_n = G \]
\[ g = t_1\,t_2\cdots t_n,\qquad t_k\in T_k\quad\text{uniquely.} \]
With indices $r_k = [G_k:G_{k-1}]$: $|G| = r_1 r_2\cdots r_n$ --- a
\alert{mixed-radix numeral}.

\vspace{4pt}
\small $S_1<S_2<\dots<S_n$: radices $2,3,\dots,n$ = the factorial number
system. The conceptual core of \alert{Schreier--Sims}; our chains are
chosen \emph{syntactically}.
\end{column}
\begin{column}{0.47\textwidth}
\centering
\begin{tikzpicture}[x=11mm,y=11mm]
  % odometer wheels
  \foreach \r/\d/\lab [count=\i] in {2/1/T_1,3/2/T_2,4/0/T_3,5/3/T_4} {
    \pgfmathsetmacro{\x}{\i*1.25}
    \draw[accent,fill=soft,rounded corners=2pt] (\x-0.5,-0.2) rectangle (\x+0.5,2.6);
    \foreach \k in {1,...,\r} {
      \pgfmathsetmacro{\yy}{2.6-\k*2.8/(\r+1)}
      \pgfmathtruncatemacro{\kk}{\k-1}
      \ifnum\kk=\d
        \node[fill=newC,text=white,font=\scriptsize,inner sep=1.5pt,rounded corners=1pt] at (\x,\yy) {\kk};
      \else
        \node[font=\scriptsize,gray] at (\x,\yy) {\kk};
      \fi
    }
    \node[font=\small] at (\x,-0.55) {$\lab$};
    \node[font=\scriptsize,gray] at (\x,2.85) {$r_\i{=}\r$};
  }
  \node[font=\small,align=center] at (3.1,-1.3) {one digit per level:\; $5!=120$ elements};
\end{tikzpicture}
\end{column}
\end{columns}
\end{frame}

\begin{frame}{Reidemeister--Schreier: certificates, not search \InPaper}
\begin{columns}[c]
\begin{column}{0.50\textwidth}
Monoid presentation $\langle X\mid R\rangle$: the quotient of $X^*$ by the
least congruence $\approx$ containing $R$; the \alert{word problem} decides
$\approx$. Interpretation $X\to G$ extends to $X^*\to G$; \emph{soundness}
descends it to $X^*/{\approx}\to G$; injective = \emph{complete}.

\vspace{6pt}
Given $H\le G=\langle Y\mid S\rangle$, a \alert{coset table} pushes a
generator $y$ past a representative:
\[ \bar t\,y \;=\; h\cdot\overline{\tau(t,y)},\qquad h\in H. \]
That is exactly \aid{ract}, and the equation is \aid{ract-sound}.
\end{column}
\begin{column}{0.47\textwidth}
\centering
\begin{tikzpicture}[node distance=6mm,
  bx/.style={draw,rounded corners=2pt,align=center,font=\small,minimum height=9mm,text width=30mm},
  >={Stealth[length=2mm]}]
  \node[bx,draw=gray,fill=gray!10] (tc) {Todd--Coxeter / by hand\\{\scriptsize search, unverified}};
  \node[bx,draw=accent,fill=soft,below=of tc] (tab) {claimed coset table\\{\scriptsize finite data}};
  \node[bx,draw=paperC,fill=softpaper,below=of tab] (chk) {Agda checks\\{\scriptsize 5 equations}};
  \node[bx,draw=paperC,fill=paperC!20,below=of chk] (nf) {normal form one\\level up, proved once};
  \draw[->,thick] (tc) -- (tab);
  \draw[->,thick] (tab) -- (chk);
  \draw[->,thick,paperC] (chk) -- node[right,font=\scriptsize]{\aid{nf-transp}} (nf);
\end{tikzpicture}

\vspace{4pt}
{\small The certificate is separated from the search that found it
(pioneered in Agda by Bian--Selinger; here a reusable module).}
\end{column}
\end{columns}
\end{frame}

\begin{frame}{For circuits the chain is handed to us \InPaper}
\begin{columns}[c]
\begin{column}{0.48\textwidth}
The shift $c\mapsto c\;{\uparrow}$ embeds
$\mathsf{Cir}\,k \le \mathsf{Cir}\,(k{+}1)$ as the circuits that do not touch
the bottom wire. So:
\begin{itemize}
  \item each \alert{wire} adds one digit, ranging over a transversal:
        staircases for $S_n$; Selinger's staircase representatives for
        Clifford groups; their qutrit analogues;
  \item a normal form for the whole group is a \alert{tower} ---
        \aid{CosetTower};
  \item where the chain \alert{branches} (two subgroups over a common base, as
        for Clifford+$T$), glue by the \alert{amalgamated free product} and its
        alternating normal form.
\end{itemize}
\end{column}
\begin{column}{0.49\textwidth}
\centering
\begin{tikzpicture}[x=6mm,y=6mm,>={Stealth[length=1.8mm]}]
  % nested tower: largest first, so the smaller ones sit on top
  \foreach \k in {4,3,2,1,0} {
    \pgfmathtruncatemacro{\shade}{6+(4-\k)*7}
    \draw[accent,fill=accent!\shade,rounded corners=2pt]
      (-\k*0.5-0.7,0) rectangle (\k*0.5+0.7,\k*1.0+0.8);
    \node[font=\scriptsize] at (0,\k*1.0+0.45) {$\mathsf{Cir}\,\k$};
  }
  \node[font=\scriptsize,align=center] at (0,-0.6) {tower: nested};
  % branching
  \begin{scope}[xshift=44mm,yshift=4mm]
    \node[draw=paperC,fill=softpaper,rounded corners=2pt,font=\scriptsize] (M) at (0,0) {$M$};
    \node[draw=paperC,fill=softpaper,rounded corners=2pt,font=\scriptsize] (A) at (-1.5,2.4) {$\langle T,M\rangle$};
    \node[draw=paperC,fill=softpaper,rounded corners=2pt,font=\scriptsize] (B) at (1.5,2.4) {$\langle H,M\rangle$};
    \node[draw=newC,fill=softnew,rounded corners=2pt,font=\scriptsize] (P) at (0,4.6) {$\langle T,M\rangle *_M \langle H,M\rangle$};
    \draw[->] (M) -- (A); \draw[->] (M) -- (B);
    \draw[->] (A) -- (P); \draw[->] (B) -- (P);
    \node[font=\scriptsize,align=center] at (0,-1.3) {branching: amalgam};
  \end{scope}
\end{tikzpicture}
\end{column}
\end{columns}
\end{frame}
